\documentclass[11pt]{amsart}
\usepackage[T1]{fontenc}
\usepackage[utf8]{inputenc}
\usepackage{amssymb}
\usepackage[margin=1.05in]{geometry}
\usepackage[hidelinks]{hyperref}
\usepackage{booktabs}
\usepackage{array}
\usepackage{tikz-cd}
\usepackage{datetime}
\numberwithin{equation}{section}
\newtheorem{theorem}{Theorem}[section]
\newtheorem{proposition}[theorem]{Proposition}
\newtheorem{lemma}[theorem]{Lemma}
\newtheorem{corollary}[theorem]{Corollary}
\theoremstyle{definition}
\newtheorem{definition}[theorem]{Definition}
\newtheorem{axiom}[theorem]{Axiom}
\theoremstyle{remark}
\newtheorem{remark}[theorem]{Remark}
\newcommand{\code}[1]{\texttt{#1}}
\newcommand{\Cat}{\mathsf{Cat}}
\newcommand{\Map}{\mathrm{Map}}
\newcommand{\Fun}{\mathrm{Fun}}
\newcommand{\Ob}{\mathrm{Ob}}
\newcommand{\Mor}{\mathrm{Mor}}
\newcommand{\Hom}{\mathrm{Hom}}
\newcommand{\Iso}{\mathrm{Iso}}
\newcommand{\Ch}{\mathrm{Ch}}
\newcommand{\dom}{\mathrm{dom}}
\newcommand{\cod}{\mathrm{cod}}
\newcommand{\id}{\mathrm{id}}
\newcommand{\II}{\mathbb{I}}
\newcommand{\one}{\mathbf{1}}
\newcommand{\isContr}{\mathrm{isContr}}
\newcommand{\isEquiv}{\mathrm{isEquiv}}
\newcommand{\Cocone}{\mathrm{Cocone}}
\newcommand{\tmap}{\mathrm{t}}
\AddToHook{shipout/foreground}{%
 \put(0.5\paperwidth,-\paperheight+0.55cm){\makebox[0pt]{\tiny\ttfamily
 unit/doc/\jobname.tex\ ---\ compiled \today\ \currenttime}}}
\title[Homology as a colimit]{Homology as a colimit in the Sattler--W\"arn
axioms,\\ and its formalisation in Unit}
\date{1 November 2026}
\begin{document}
\sloppy
\maketitle

\section{Introduction}

In homotopy type theory the cohomology of a type $X$ with coefficients in an
abelian group $A$ is simple to define:
\[ H^n(X;A) = \| X \to K(A,n) \|_0 . \]
The function type $X \to K(A,n) = \prod_{x:X} K(A,n)$ is the limit over $X$ of
the constant diagram at $K(A,n)$. Limits of groups, spectra or chain complexes
are computed on the underlying types, so $\Pi$ suffices.

Homology is the dual notion: the colimit over $X$ of the constant diagram at
the coefficients, taken in a category $D$ of abelian objects (spectra, or chain
complexes up to quasi-isomorphism). This colimit is not $\Sigma$, because
forgetful functors do not preserve colimits. Type theory has native colimits
in one category only, the universe. Homology therefore needs colimits indexed
by types in a different category. D.~W\"arn suggested that a theory of higher
categories inside type theory could provide them.

This note tests the suggestion in the axiomatic setting of Sattler and W\"arn,
which we formalised earlier in Unit (\code{unit/doc/hcat-segal.pdf}). For a
constant diagram the colimit is the \emph{copower} $X\otimes A$. We define
homology as $X\otimes A$ and prove the Eilenberg--Steenrod axioms at the level
of $D$:
\begin{itemize}
\item homotopy invariance;
\item dimension, $\one\otimes A\cong A$;
\item additivity;
\item excision: pushouts of types go to pushouts in $D$;
\end{itemize}
as well as a Fubini formula. With a zero object and stability we then define
the homology groups, prove the Mayer--Vietoris sequence (in part) and the
suspension isomorphism, and compute $\tilde H_*(S^n)$ (\S\ref{sec:groups}).
All of the proofs of the axioms are short consequences of the
universal property except excision. Excision needs the pentagon for
composition inside $D$. In this setting that is a coherence of level~3 of
$\Cat$ (\S\ref{sec:coh}).

\section{The setting}

We work in dependent type theory with univalence, as in \code{hcat-segal}.
$\Cat$ is a wild category: a type of categories with types $\Map(C,D)$,
composition and identities, and its coherences given as axioms level by level.
For a category $C$:
\begin{itemize}
\item $\Ob C=\Map(\one,C)$;
\item $\Mor C=\Map(\II,C)$;
\item $\Hom_C(x,y)$ consists of the morphisms with endpoint paths to $x$ and $y$.
\end{itemize}
Composition $gf$ of $g:\Hom(x,y)$ and $f:\Hom(y,z)$ (written in diagrammatic
order) is the long edge of the centre of the contractible type of 2-simplices
with spine $(g,f)$; this uses the Segal axiom. Associativity
\[ \alpha(f,g,h) : (fg)h = f(gh) \]
comes from Segal at $n=3$ and uses the pentagon of $\Cat$. The unit laws are
similar. Every category is univalent (the Rezk axiom). In particular, a functor
from a groupoid $X$ into $D$ is a function $X\to\Ob D$.

\section{Copowers}

\begin{definition}
Let $D$ be a category. A \emph{copower} of $A:\Ob D$ by a type $X$ is an object
$X\otimes A$ with maps $\eta_x : A\to X\otimes A$ for $x:X$, such that for
every $T$
\[ \mathrm{restrict} : \Hom(X\otimes A,T)\to (X\to \Hom(A,T)),\qquad
   g\mapsto (x\mapsto \eta_x g) \]
is an equivalence. We write $\mathrm{lift}$ for its inverse.
\end{definition}

$X\otimes A$ is the colimit of the constant $X$-diagram at $A$, and we read it as
the homology of $X$ with coefficients $A$. In the intended model $D$ is the
category of spectra and $X\otimes A=\Sigma^\infty X_+\wedge A$, so that
$\pi_n(X\otimes HZ)=H_n(X;Z)$.

\begin{remark}[Size]
The objects of a category form a type in a universe $U_l$. We assume copowers
only by types in a strictly smaller universe $U_k$, $k<l$. With copowers by
all types of $U_l$, Freyd's argument would make $D$ a preorder: take
$X=\Sigma_C \Hom(C,T)$ and compare the sizes of
$\Hom(X\otimes A,T)$ and $X\to\Hom(A,T)$.
\end{remark}

Two consequences of the universal property are used everywhere.
\begin{itemize}
\item \emph{$\beta$-law}: $\eta_x\,\mathrm{lift}(h)=h(x)$.
\item \emph{Extensionality}: two maps out of $X\otimes A$ that agree after every
  $\eta_x$ are equal, since restrict is injective and there is function
  extensionality.
\end{itemize}

\section{The axioms without 2-cells}

For $f:X\to Y$ let $\tmap f = \mathrm{lift}(x\mapsto \eta_{f(x)}) : X\otimes A\to
Y\otimes A$. Then $\eta_x\,\tmap f=\eta_{f(x)}$, and by associativity
\begin{equation}\label{eq:rt}
\eta_x(\tmap f\, g)=\eta_{f(x)}\,g .
\end{equation}

\begin{proposition}[Functoriality]
$\tmap(\id)=\id$, $\tmap(g\circ f)=\tmap f\,\tmap g$, and homotopic maps have
equal images.
\end{proposition}
\begin{proof}
By extensionality. After $\eta_x$ both sides of each equation become
$\eta_x$, $\eta_{g(f(x))}$, and $\eta_{f(x)}=\eta_{f'(x)}$ respectively, by the
$\beta$-law and \eqref{eq:rt}.
\end{proof}

\begin{theorem}
\begin{enumerate}
\item (Homotopy invariance) If $f:X\to Y$ is an equivalence, then $\tmap f$ is
  invertible.
\item (Dimension) $\eta_\star : A\to \one\otimes A$ is invertible.
\item $0\otimes A$ is initial.
\item (Additivity) $(X+Y)\otimes A$, with $\tmap\,\mathrm{inl}$ and
  $\tmap\,\mathrm{inr}$, is a coproduct of $X\otimes A$ and $Y\otimes A$.
\item (Fubini) $(X\times Y)\otimes A\cong X\otimes(Y\otimes A)$.
\end{enumerate}
\end{theorem}
\begin{proof}
(1) The inverse is $\tmap f^{-1}$, by functoriality.

(2) The inverse is $r=\mathrm{lift}(\_\mapsto\id_A)$. We have $\eta_\star r=\id$
by the $\beta$-law. For $r\eta_\star=\id$ use extensionality:
$\eta_\star(r\eta_\star)=(\eta_\star r)\eta_\star=\eta_\star$.

(3) $\Hom(0\otimes A,T)\simeq(0\to\Hom(A,T))$ is contractible.

(4) The inverse of $g\mapsto(\tmap\,\mathrm{inl}\,g,\tmap\,\mathrm{inr}\,g)$
sends $(h_1,h_2)$ to the lift of the family
$\mathrm{inl}\,x\mapsto\eta_x h_1$, $\mathrm{inr}\,y\mapsto \eta_y h_2$. Both
composites are identities by extensionality and \eqref{eq:rt}.

(5) The map $(X\times Y)\otimes A\to X\otimes(Y\otimes A)$ is the lift of
$(x,y)\mapsto \eta_y\eta_x$. Its inverse is the lift of
$x\mapsto\mathrm{lift}(y\mapsto\eta_{(x,y)})$. Both composites are identities by
extensionality, applied twice for the second one, together with associativity
and the $\beta$-law.
\end{proof}

Each proof uses only associativity and the unit laws of $D$. So all of them
hold at level~2 of $\Cat$.

\section{Excision}

There are no higher inductive types here, so a pushout of types is
characterised by its universal property.

\begin{definition}
A commuting square of types
\[\begin{tikzcd} C\ar[r,"i"]\ar[d,"j"'] & B\ar[d,"u"]\\ E\ar[r,"v"'] & X
\end{tikzcd}\qquad h(c): u(i(c))=v(j(c))\]
is a \emph{pushout} if for every type $T$ (in the universe of the hom types of
$D$) the map
$f\mapsto (f\circ u, f\circ v, \mathrm{ap}_f\circ h)$ from $X\to T$ to the type
$\Cocone(T)$ of triples $(b,e,\theta)$, with
$\theta(c):b(i(c))=e(j(c))$, is an equivalence.

A square $i'u'=j'v'$ in $D$, with 2-cell $s$, is a \emph{pushout} if for every
$T$ the map
\[ g\mapsto \big(u'g,\ v'g,\ \alpha^{-1}\cdot (s\,g)\cdot\alpha\big) :
  \Hom(P,T)\to \Cocone_D(T) \]
is an equivalence. Here $\Cocone_D(T)$ consists of triples $(b,e,p)$ with
$p:i'b=j'e$, and the path part is
$i'(u'g)=(i'u')g=(j'v')g=j'(v'g)$.
\end{definition}

\begin{theorem}[Excision]
Assume the pentagon for composition in $D$. Then a pushout square of types is
sent to a pushout square in $D$:
$\tmap i,\tmap j,\tmap u,\tmap v$ with the 2-cell $\sigma$ defined by
$\eta_c\,\sigma = $ the path
$\eta_c(\tmap i\,\tmap u)=\eta_{u(i(c))}=\eta_{v(j(c))}=\eta_c(\tmap j\,\tmap v)$
built from \eqref{eq:rt} and $h(c)$.
\end{theorem}
\begin{proof}
Let $T:\Ob D$ and $H=\Hom(A,T)$. Define
$\psi:\Cocone_D(T)\to\Cocone(H)$ by
\[ (b,e,p)\mapsto(\mathrm{restrict}\,b,\ \mathrm{restrict}\,e,\ \theta_p),\qquad
\theta_p(c) = \eta_c\,p \text{ conjugated by \eqref{eq:rt} at } i \text{ and } j. \]
$\psi$ is an equivalence: restrict is an equivalence on both legs, and on
paths $p\mapsto(c\mapsto\eta_c\,p)$ is an equivalence, being $\mathrm{ap}$ of
restrict followed by function extensionality.

The square of maps
\[\begin{tikzcd}[column sep=large]
\Hom(X\otimes A,T)\ar[r,"\text{comparison}_D"]\ar[d,"\mathrm{restrict}"'] &
\Cocone_D(T)\ar[d,"\psi"]\\ (X\to H)\ar[r,"\text{comparison}"'] & \Cocone(H)
\end{tikzcd}\]
commutes up to homotopy. On the legs this is \eqref{eq:rt}. On the path part
we compare, for each $c$, two paths in $H$. Both are composites of:
\begin{itemize}
\item the associators at $(\eta_c,\tmap i,\tmap u,g)$ and at
  $(\eta_c,\tmap j,\tmap v,g)$;
\item instances of \eqref{eq:rt} whiskered by $g$;
\item the path $\mathrm{ap}_\eta(h(c))$.
\end{itemize}
Each instance of \eqref{eq:rt} contains an associator, and composing it with
the associator of the comparison map rebrackets the four-fold composite
$\eta_c\,(\tmap i)\,(\tmap u)\,g$. Two such rebracketings agree exactly by the
pentagon at $(\eta_c,\tmap i,\tmap u,g)$, and similarly at $j$. After that the
equation is a groupoid identity.

The left and bottom maps are equivalences, by the copower and by the pushout
of types at $H$. Hence the top map is one.
\end{proof}

\section{The coherence needed}\label{sec:coh}

The pentagon for composition inside a category $C$ is not an axiom: composition
and $\alpha$ are defined from contractible spaces of simplices. To prove it
from the Segal condition at $n=4$ one compares the five associators of a
4-simplex $T$ through its faces. That requires relating
\begin{itemize}
\item composites of the form $T\circ S\circ s$, where $S$ is a face and $s$ a
  2-simplex of it, with
\item the action of $T$ on the composites of $S$.
\end{itemize}
This is a coherence for composition in $\Cat$ one level up: \emph{internal}
level~2 (the pentagon of $C$) needs \emph{external} level~3. It is the same
shift as one level lower, where associativity in $C$ needs the pentagon of
$\Cat$. It is also the effect described by Sattler: an $n$-coherent copy of
$\Cat$ inside $\Cat$ needs $(n+1)$-coherence.

The level~3 coherence of $\Cat$ is Stasheff's associahedron $K_5$. Take five
maps $d,e,h,g,f$ and write $a(x,y,z):(xy)z=x(yz)$. There are two extreme paths
from $(((fg)h)e)d$ to $f(g(h(ed)))$:
\begin{itemize}
\item $L$, the three-step path $a((fg)h,e,d)\cdot a(fg,h,ed)\cdot a(f,g,h(ed))$;
\item $R$, the six-step path through $((f(gh))e)d$.
\end{itemize}
There are two ways to paste the faces of $K_5$ into a 2-path $L=R$:
\begin{itemize}
\item route~1 uses the pentagons at $(f,g,h,ed)$, $(f,gh,e,d)$ and $(g,h,e,d)$
  (the last whiskered by $f$), and the naturality squares of $a(-,e,d)$ and
  $a(f,-,d)$;
\item route~2 uses the pentagons at $(fg,h,e,d)$, $(f,g,he,d)$ and
  $(f,g,h,e)$ (the last whiskered by $d$), and the naturality square of
  $a(f,g,-)$.
\end{itemize}
The axiom says the two routes agree.

In the proof of the internal pentagon, K5 is used where the action of functors
on morphisms (with their endpoints) is compared for three functors at once.
At $(F,G,K,m,\II_0)$, the endpoint 2-cells of the two sides follow route~1 and
route~2.

If composition in $\Cat$ were strictly associative and unital, the pentagons
and $K_5$ would be trivial, and the internal pentagon would follow from Segal
at $n=4$ by routine path algebra. But it would no longer be visible which
coherence the argument uses.

\section{Homology groups}\label{sec:groups}

\subsection*{Stability}
Assume a zero object $0$ (both $\Hom(0,T)$ and $\Hom(B,0)$ contractible), write
$0$ also for the zero maps, and assume suspensions: for each $B$ a pushout
square $0\leftarrow B\to 0$ with apex $\Sigma B$ and 2-cell $m_B$. For
\emph{stability} we take the representable form of ``pushout squares are
pullbacks'':
\begin{axiom}[Stability]
For every pushout square $i u=j v$ in $D$ (with $i:C\to B$, $j:C\to E$) and every
$S$, the map $c\mapsto(ci,\ cj,\ \alpha\cdot c\,s\cdot\alpha^{-1})$ from
$\Hom(S,C)$ to the type of triples $(b,e,p)$, $p:bu=ev$, is an equivalence.
\end{axiom}
It needs no pullbacks in $D$, and it is exactly what the arguments below use.
Of ``$\Sigma$ is an equivalence'' it gives full faithfulness; essential
surjectivity (desuspensions) is used only for $\tilde H_k(S^n)$ with $k\le n$.

\begin{lemma}[Loops]\label{lem:loops}
$L:\Hom(\Sigma S,T)\simeq\Omega(\Hom(S,T),0)$, natural in $T$: for
$f:T\to T'$, $L(gf)$ is $L(g)f$ conjugated by the path $0f=0$, and $L(0)=\mathrm{refl}$.
With stability also $\sigma:\Hom(B,E)\simeq\Omega(\Hom(B,\Sigma E),0)$.
\end{lemma}
\begin{proof}
$\Hom(\Sigma S,T)$ is the type of cocones $(b,e,p)$ with $b,e:0\to T$; as
$\Hom(0,T)$ is contractible this is the type of loops at $0$. Naturality: both
sides are conjugates of $m_S(gf)$, and the two conjugating paths agree by the
pentagon at $(0,0,g,f)$ and naturality of $\alpha$; this is the second use of
the pentagon of $D$. $L(0)=\mathrm{refl}$ since $0:\Sigma S\to T$ factors
through $0$, and $\Hom(S,0)$ is contractible. $\sigma$ is stability for
$0\leftarrow E\to 0$.
\end{proof}

\begin{definition}
Let $\one:\Ob D$ and $S^n_D=\Sigma^n\one$. Then
$H_n(X;A)=\pi_0\Hom(S^n_D,X\otimes A)$, and reduced homology is
$\tilde H_n(X;A)=\pi_0\,\mathrm{Red}(S^n_D,X)$, where $\mathrm{Red}(S,X)$ is the
fibre over $0$ of $c\mapsto c\,\tmap(!)$, $\Hom(S,X\otimes A)\to\Hom(S,\one\otimes A)$.
\end{definition}
$\mathrm{Red}(S,X)$ is $\Hom(S,-)$ of the fibre of $X\otimes A\to \one\otimes A$, so
no cofibres in $D$ are needed.

\begin{theorem}[Groups]
$H_{n+1}(X;A)$ is a group, through $\pi_0 L$ from $\pi_1(\Hom(S^n_D,X\otimes A),0)$,
and it is abelian for $n\ge1$ (Eckmann--Hilton through $L$). With stability
$H_0(X;A)$ is a group (through $\sigma$) and all these groups are abelian.
\end{theorem}

\subsection*{Mayer--Vietoris}
\begin{theorem}
Let $iu=jv$ be a pushout square in $D$ and $S:\Ob D$; write $i_*$ etc.\ for
postcomposition on $\pi_0\Hom(S,-)$.
\begin{enumerate}
\item $u_*i_*=v_*j_*$, and if $u_*\beta=v_*\varepsilon$ then merely
  $\beta=i_*\gamma$ and $\varepsilon=j_*\gamma$ for some $\gamma$.
\item There is $\delta:\pi_0\Hom(\Sigma S,P)\to\pi_0\Hom(S,C)$ with
  $i_*\delta=0$, $j_*\delta=0$, and every $\gamma$ with $i_*\gamma=0=j_*\gamma$ is
  merely $\delta$ of an element.
\end{enumerate}
\end{theorem}
\begin{proof}
By stability $\Hom(S,C)$ is the type of triples $(b,e,p:bu=ev)$. (1): a path
$bu=ev$ gives the triple, hence $c$. (2): $\delta g$ corresponds to the triple
$(0,0,p_g)$ where $p_g$ is $L(g)$ conjugated by $0u=0$ and $0v=0$; conversely
paths $ci=0$, $cj=0$ move the triple of $c$ to one of this form.
\end{proof}
With excision this gives, for a pushout of types,
$H_{n+1}(X)\xrightarrow{\delta}H_n(C)\to H_n(B)\times H_n(E)\to H_n(X)$, exact at
$H_n(C)$ and in the middle.

\begin{theorem}[Positive-degree exact sequence]
For every $r\geq2$, a single structural connecting homomorphism gives
\[
\begin{aligned}
 H_r(B)\times H_r(E)&\xrightarrow{u_*-v_*}H_r(X)
 \xrightarrow{\partial}H_{r-1}(C)\\
 &\xrightarrow{(i_*,j_*)}H_{r-1}(B)\times H_{r-1}(E)
 \xrightarrow{u_*-v_*}H_{r-1}(X).
\end{aligned}
\]
Every arrow is a homomorphism of the specified abelian groups and the
sequence is exact at its three interior terms. Consecutive pieces agree.
\end{theorem}
\begin{proof}
For any pointed map, the equivalence between loops of its fibre and
the fibre of its loop map retains the projection. Taking components
therefore gives both homomorphisms and exactness for the same maps.
Apply this first to the projection of the representable pullback onto
the product of its legs, whose fibre is the loop space of the apex.
Apply it again to the boundary: its fibre is the loop space of that
product, over the difference of loop maps. Looping difference gives
ordinary group difference on $\pi_2$.

The double-loop equivalence obtained from $L$ identifies these with
the existing homology groups and maps. Its naturality uses $L$'s
already proved naturality and one general lemma: homotopic pointed
maps into a type with commuting loops induce equal loop maps,
regardless of their chosen pointing paths. Conjugation by a loop
is trivial, so no new compatibility of those paths is needed.
Exactness at the product is the preceding pullback argument and group
cancellation. Difference is additive because the target is abelian.
\end{proof}
This is the positive tail ending at $H_1(X)$, with no exactness claim
at that final term. Its extension through degree zero and categorical
naturality remain open. The structural boundary need not be compared
with the earlier direct formula to obtain these results.

\subsection*{Suspension and spheres}
A \emph{suspension} of a type $X$ is a pushout $\Sigma X$ of $\one\leftarrow X\to\one$,
with points $N,S$; it is given by its universal property.
\begin{theorem}[Suspension]
$\mathrm{Red}(\Sigma S,\Sigma X)\simeq\mathrm{Red}(S,X)$; hence
$\tilde H_{n+1}(\Sigma X;A)\cong\tilde H_n(X;A)$.
\end{theorem}
\begin{proof}
Let $P=\Sigma X\otimes A$, $t_N,t_S:\one\otimes A\to P$ and $r:P\to\one\otimes A$
the maps of $N$, $S$ and $\Sigma X\to\one$, and $F=\mathrm{Red}(S,\Sigma X)$,
the fibre of $r_*$, pointed at $0$. We show that both sides are $\Omega F$.

By excision and stability, $\Hom(S,X\otimes A)$ is the type of triples
$(a,b,p:at_N=bt_S)$, the first leg being $c\mapsto c\,\tmap(!)$. So
$\mathrm{Red}(S,X)$ is the type of pairs $(b,\ 0t_N=bt_S)$, i.e.\ the fibre of
$(t_S)_*$ over $0$. Since $t_Sr=\tmap(\one\to\one)=\id$, $(t_S)_*$ has the retraction
$r_*$. For any such pair of maps the fibre of $s$ over $s(y_0)$ is
$\Omega$ of the fibre of $r$ (the total space of the composite fibration is
contractible). Moving base points along zero paths gives $\Omega F$.

On the other side, change the base of $\mathrm{Red}(\Sigma S,\Sigma X)$ along $L$:
$gr=0$ iff $L(gr)=L(0)$, i.e.\ by Lemma~\ref{lem:loops} iff $L(g)\,r$ conjugated
is $\mathrm{refl}$. So it is the fibre over $\mathrm{refl}$ of $\Omega r_*$, which
is $\Omega F$.
\end{proof}

Let $S^0=\Sigma\emptyset$ and $S^{n+1}=\Sigma S^n$. Since $\emptyset\otimes A$ is
initial, $\Hom(S,\emptyset\otimes A)$ is contractible, so
$\mathrm{Red}(S,\emptyset)\simeq\Omega\Hom(S,\one\otimes A)\simeq\Hom(\Sigma S,\one\otimes A)$.
\begin{corollary}[Spheres]
$\tilde H_{n+j+1}(S^n;A)\cong\pi_0\Hom(S^{j+1}_D,A)=H_{j+1}(\one;A)$. With
desuspensions ($\Sigma\Sigma^{-1}B=B$), $\tilde H_{n+j}(S^n;A)\cong\pi_0\Hom(S^j_D,A)$
for all $j$, in particular $\tilde H_n(S^n;A)\cong\pi_0\Hom(\one,A)=H_0(\one;A)$,
and $\tilde H_k(S^{k+m};A)\cong\pi_0\Hom(\Sigma^{-m}\one,A)$.
\end{corollary}
\begin{proof}
The suspension theorem $n$ times, then $\emptyset$, and the dimension axiom
$\one\otimes A\cong A$. With desuspensions write $B=\Sigma\Sigma^{-1}B$ first.
\end{proof}
These are bijections of sets; compatibility of the suspension and sphere
isomorphisms with the specified group structures remains unproved.

\section{Formalisation}

The development is in \code{unit/sag/hcat}, continuing the port of
\code{hcat-segal}. Segments are nested by level: \code{Theory} $\supset$
\code{Wild} $\supset$ \code{Structure} $\supset$ \code{Coh2}. Inside
\code{Coh2} there are two new segments.
\begin{itemize}
\item \code{Coh3}, with parameter \code{k5 : route1 = route2}; the routes are
  defined in \code{CatK5}.
\item \code{Homology}, with parameters: a level $k<l$, the category
  \code{D}, \code{tensor}, \code{eta}, and \code{tensor-up} (restrict is an
  equivalence). Inside it, the sub-segment \code{WithPentagon} has the
  parameter \code{hpent}, the pentagon of $D$ (stated in
  \code{HcompPentagon}). Inside that, \code{Stable} has the zero object,
  suspensions, \code{stable} and $\one$; inside it \code{WithSpheres} has the
  suspensions of types and the desuspensions.
\end{itemize}

\begin{center}
\begin{tabular}{>{\raggedright}p{6.2cm}l}
\toprule
Result & File, name \\
\midrule
lift, $\beta$-law, extensionality & \code{Copower}: \code{lift}, \code{lift-beta}, \code{ext}\\
functoriality, \eqref{eq:rt} & \code{TensorMap}: \code{tmap-id}, \code{tmap-comp}, \code{restrict-tmap}\\
homotopy invariance, dimension, initial & \code{TensorInvariance}\\
additivity & \code{TensorSum}: \code{SumOf.tensor-sum}\\
Fubini & \code{Fubini}: \code{FubiniOf.fubini}\\
pushouts of types, homotopy induction & \code{TypePushout}\\
pushouts in $D$ & \code{DPushout}\\
the pentagon step & \code{ExcisionAssoc}: \code{rt-assoc}\\
excision & \code{Excision}: \code{ExcisionOf.excision}\\
$K_5$ & \code{CatK5}, parameter \code{Coh3.k5}\\
zero maps, cones & \code{Zero}, \code{Cone}\\
Lemma~\ref{lem:loops} & \code{Loops}: \code{L}, \code{L-nat}, \code{L-zero}; \code{StableLoop}\\
groups & \code{HGroups}: \code{hgrp}, \code{hgrp-ab}, \code{hgrp0}, \code{hgrp-ab-stable}\\
Mayer--Vietoris & \code{MayerVietoris}: \code{MV}, \code{MVTypes}\\
positive-degree exact sequence & \code{HomologySequence}: \code{Types}\\
suspension & \code{Suspension}: \code{SuspOf.susp-red}, \code{susp-iso}\\
spheres & \code{Spheres}: \code{sphere-above}, \code{sphere-homology}\\
\bottomrule
\end{tabular}
\end{center}

\noindent The pentagon of $D$ (\code{hpent}) supplies the
lemma \code{rt-assoc} of excision, \code{L-nat}, and the compatibility
of zero-postcomposition paths in \code{ZeroCoherence}. The type theory of
\S\ref{sec:groups} (sets of components, loop groups, Eckmann--Hilton, fibres) is in
\code{Pi0}, \code{EckmannHilton}, \code{FibreAlg}, \code{Conj}, partly linked from
\code{sag/em}, with the further fibre and difference lemmas described
above. The tree has 169 modules and 3978 declarations; a fresh full
check with unfolding diagnostics takes 17.44\,s with no failures.
Each new file checks in under 1\,s. The homology sequence does not
depend on \code{Coh3}.

There are no new axioms beyond:
\begin{itemize}
\item the parameters \code{univ}, \code{images} and \code{surjections} of
  \code{Theory};
\item the axioms of $\Cat$ in \code{Wild}, \code{Structure}, \code{Coh2} and
  \code{Coh3};
\item the parameters of \code{Homology};
\item \code{hpent};
\item the parameters of \code{Stable}: \code{zobj}, \code{zero-init},
  \code{zero-term}, \code{susp}, \code{merid}, \code{susp-po}, \code{stable},
  \code{unitD};
\item the parameters of \code{WithSpheres}: \code{tsusp}, \code{tnorth},
  \code{tsouth}, \code{tmerid}, \code{tsusp-po}, \code{desusp},
  \code{desusp-eq}.
\end{itemize}

\noindent The formal proofs include the path, fibre and quotient lemmas
used in these arguments. The positive exact sequence is assembled from
those reusable lemmas in \code{HomologySequence}; it adds no axioms.

\subsection*{Open}
\begin{itemize}
\item Deriving \code{hpent} from \code{k5}. The plan is in
  \code{sag/hcat/DESIGN-HOMOLOGY.md}; one lemma of the needed calculus of
  paths between morphisms with endpoints (\code{HomEqCalc}) is done.
\item Degree-zero homomorphisms and the extension of the structural
  exact sequence through degree zero; categorical naturality of that
  sequence; the group form of the suspension isomorphism.
\end{itemize}

\end{document}
